\section{Stand der Technik}

\subsection{Marktrecherche}
Im Internet sind viele Werkzeughalter zu finden. Diese sind jedoch auf bestimmte Werkzeuge abgestimmt und nicht auf vereinzelte Maschinenelemente. Als Beispiel ein Wandhalter für Werkzeug von \textit{Amazon} \cite{Amazon2026}. Zu sehen auf Abbildung \ref{fig:amazonHalter}:

\begin{figure}[H]
	\centering
	\includegraphics[width=0.5\linewidth]{amazonHalter.png}
	\caption{Kommerzieller Halter für Werkzeug}
	\label{fig:amazonHalter}
\end{figure}

\subsection{Lösungsansätze}
\textcolor{red}{Skizziert idealerweise mehrere Konzeptvarianten, bevor bewertet wird – z. B. Steckboard/Lochraster, Einzelsockel pro Element mit eingeprägter Bezeichnung, Klemm-/Rastlösung, geneigte Präsentationsfläche.}

\subsection{Konzeptvergleich}
\textcolor{red}{Eine formale Bewertungsmatrix mit gewichteten Kriterien wäre hier stark – die Kriterien liefert eigentlich schon eure Zielsetzung: Standsicherheit, Materialverbrauch/Druckzeit, Übersichtlichkeit/Funktionssichtbarkeit, Informationsvermittlung. -> Wie in MeKo damals}